\scanpage{9}
\mspage{7}
(2.2.2), et on ne se bornera pas \deleted{\unsure{et|à}} à étudier sa restriction à \(I\). \deleted{\unsure{C'est|le}} On verra \deleted{\unsure{parfois|toutefois}} que \inserted{parfois} le fait qu'une représentation de \(I\) se prolonge à \(\pi\) impose à cette représentation des restrictions draconiennes.

2.3.

2.4. Rappelons la structure de \(I\). Soit \(p\) l'exposant caractéristique de \(k\), \deleted{Alors} et soit \(P\) \deleted{le} un \deleted{plus grand} sous-\(p\)-groupe de Silow de \(I\) (\(P=1\) si \(p=1\) i.e. \(k\) de car. nulle). Alors \(P\) \uline{est invariant}, i.e. on a une suite exacte
\[
1\mathrel{\unsure{-|\to}}P\longrightarrow I\mathrel{-} I_t\mathrel{-}1,\tag*{(2.4.1)}
\]
où \(I_t\) \inserted{(partie ``tame'' du groupe d'inertie \(I\))} est un groupe profini premier à \(p\), d'autre part on a un isomorphisme canonique
\[
I_t\simeq\prod_{\ell\ne p}\mathbb Z_\ell(1)\tag*{(2.4.2)}
\]

Ce dévissage est obtenu ainsi. Le groupe \(P\) est le groupe de Galois \deleted{de}
\[
P=\operatorname{Gal}(\overline K/\overline K_t),\tag*{(2.4.1)}
\]
où \(\overline K_t\) est la plus \deleted{\unsure{sous|grande}} grande sous-extension \uline{modérément ramifiée} de \(\overline K/K\) \inserted{(rel. à \(V\))}, ou encore de \(\overline K/\widetilde K\) (rel. à \(\widetilde V\)). Celle-ci s'obtient, en termes
% archive sheet number: 8
